\section{Instantiation}\label{sec:offline-instantiation}

We now instantiate the generic construction of
\Cref{sec:offline-construction} with concrete cryptographic building
blocks. We use BBS+ signatures to build the anonymous credential that binds
a user's pseudonym to their secret key (\Cref{sec:BBS}), and
Pointcheval-Sanders signatures, extended to support blind signing over
committed inputs, to let the central bank sign tokens at withdrawal without
learning their content (\Cref{sec:ps-sign}). Both are pairing based
multi-message signature schemes, which lets us build efficient sigma
protocols directly over their verification equations rather than relying on
general purpose zero-knowledge proofs.

\paragraph{Notation.} Let $[n]$ denote the integer interval $[0..n-1]$. Let
$\group{1}$, $\group{2}$, and $\group{T}$ be three cyclic groups of large
prime order $p$ that admit a non-degenerate bilinear pairing $\pairing
\colon \group{1} \times \group{2} \rightarrow \group{T}$. Let $\hash_p
\colon \{0,1\}^* \rightarrow \Zp$ and $\hash_{\group{1}} \colon \{0,1\}^*
\rightarrow \group{1}$ be two cryptographic hash functions. Let $\gen$ and
$\ggen$ denote generators of $\group{1}$ and $\group{2}$ respectively. We
denote elements of $\group{1}$ by lower-case letters, elements of
$\group{2}$ by lower-case letters with a tilde, and elements of $\Zp$ by
lower-case Greek letters.

For a vector $\vec{x} = (x_0, \ldots, x_{n-1}) \in \Zp^n$ and $y \in \Zp$,
and $m < n - 1$, we write $\vec{x}[{:}m]$ for the sub-vector $(x_0, \ldots,
x_{m-1})$ and $\vec{x} \| y$ for $(x_0, \ldots, x_{n-1}, y)$, with the
analogous notation for vectors in $\group{1}$ and $\group{2}$. For
$\vec{x}, \vec{y} \in \Zp^n$, we write $\vec{x} \cdot \vec{y}$ for the
inner product $\sum_{i=0}^{n-1} x_i y_i$. For $\vec{\gen} = (\gen_0,
\ldots, \gen_{n-1}) \in \group{1}^n$, $\vec{\ggen} = (\ggen_0, \ldots,
\ggen_{n-1}) \in \group{2}^n$, and $\vec{x} = (x_0, \ldots, x_{n-1}) \in
\Zp^n$, we write $\vec{\gen}^{\vec{x}} = \prod_{i=0}^{n-1} \gen_i^{x_i}$
and $\vec{\ggen}^{\vec{x}} = \prod_{i=0}^{n-1} \ggen_i^{x_i}$.

We write $\algoproof \leftarrow \NIZKprove(\crs, \inst, \witness)$ for a
non-interactive zero-knowledge proof of knowledge that public input
$\inst$ and witness $\witness$ satisfy the accompanying relation, and
$\{0,1\} \leftarrow \NIZKverify(\crs, \inst, \algoproof)$ for its
verification, where $\crs$ is the common reference string published by
$\algo{Setup}$ (\Cref{sec:offline-construction}). Appendix \ref{sec:zkp}
gives an overview of zero-knowledge proofs.

\subsection{BBS+ signatures}\label{sec:BBS}

A BBS+ signature is a signature scheme
that lets a signer sign multiple messages at once, using the following
algorithms~\cite{C:BonBoySha04,SCN:AuSusMu06,EPRINT:CamDriLeh16}.

\begin{primitivebox}{BBS+ signatures}{bbs}

  \begin{description}

    \item[$(\sk_{\BBS}, \pk_{\BBS}) \leftarrow \BBSgen(1^\kappa, n)$ :] On input
      security parameter $\kappa$ and integer $n$, randomly select $\theta
      \in \Zp$ and a vector of $n+2$ generators $(q_0, \ldots, q_n, \hgen)
      \in \group{1}^{n+2}$, writing $\vec{q} = (q_0, \ldots, q_n)$. Compute
      $\tilde{t} = \ggen^\theta$ and output secret key $\sk_{\BBS} = \theta$
      and public key $\pk_{\BBS} = (\tilde{t}, \ggen, \vec{q}, \hgen)$.

    \item[$\phi \leftarrow \BBSsign(\sk_{\BBS}, \vec{m})$ :] On input secret key
      $\sk_{\BBS}$ and a vector $\vec{m}$ of $n$ messages, randomly select
      $(\varepsilon, s) \in \Zp$ such that $\varepsilon + \theta \neq 0$,
      compute $c = \hgen \vec{q}^{\vec{m}\|s}$, and output $\phi =
      (\varepsilon, s, e = c^{1/(\theta+\varepsilon)})$.

    \item[$\{0,1\} \leftarrow \BBSverify(\pk_{\BBS}, \vec{m}, \phi)$ :] On input
      public key $\pk_{\BBS}$, vector of messages $\vec{m}$, and signature
      $\phi$, compute $c = \vec{q}^{\vec{m}\|s}\hgen$ and return the result of
      the check given as follows:
      \begin{equation*}
        \pairing(e, \ggen^\varepsilon \tilde{t}) = \pairing(c, \ggen).
      \end{equation*}

  \end{description}

\end{primitivebox}

BBS+ signatures are existentially unforgeable under the $q$-strong
Diffie-Hellman ($q$-SDH) assumption.

\begin{assumptionbox}{$q$-Strong Diffie-Hellman ($q$-SDH)}{qsdh-bbs}

  The $q$-SDH assumption holds if, given $(\hgen, \ggen, \ggen^{\theta},
  \gen_0, \gen_0^{\theta}, \ldots, \gen_0^{\theta^{q}})$ for uniformly
  random $\theta \sample \Zp$, no PPT adversary can output a pair
  $(\varepsilon, c^{1/(\theta+\varepsilon)})$ for a group element $c$ of
  this form and any $\varepsilon \in \Zp \setminus \{-\theta\}$ with
  non-negligible probability.

\end{assumptionbox}

\paragraph{Proof of knowledge of a BBS+ signature.}\label{sec:joint-proof}
Let $\phi = (\varepsilon, s, e) \gets \BBSsign(\sk_{\BBS}, \vec{m})$. By
construction $e = (\hgen \vec{q}^{\vec{m}\|s})^{1/(\theta+\varepsilon)}$
for $\sk_{\BBS} = \theta$; setting $\inst =
\pk_{\BBS} = (\ggen, \tilde{t}, \vec{q}, \hgen)$ and $\witness = (\vec{m}, \phi)$, a
proof of knowledge of a BBS+ signature shows that $(\inst, \witness)$
satisfies the equation as follows:
\begin{equation*}
  \pairing(e, \ggen^\varepsilon \tilde{t}) = \pairing(\hgen \vec{q}^{\vec{m}\|s}, \ggen).
\end{equation*}
A prover can instead compute $u =
e^{-\iota/\lambda}$, $w = (\hgen \vec{q}^{\vec{m}\|s})^{-1/\lambda}$, and $x =
u^\varepsilon w^\iota$, for randomly selected $\iota, \lambda$, and prove that for $\inst =
(\pk_{\BBS}, u, w, x)$ and $\witness = (\varepsilon, s, \iota, \lambda, \vec{m})$ the relation as follows:
\begin{align}
  \pairing(u, \tilde{t}) \pairing(x, \ggen) = 1 \ \wedge \ x = u^\varepsilon w^\iota \ \wedge \
  \hgen^{-1} = w^\lambda \vec{q}^{\vec{m}\|s}.
  \label{eq:BBSrel}
\end{align}
This holds~\cite{EC:TesZhu23b}. A prover $\algo{Prv}$ for this relation follows a sigma protocol made
non-interactive via the Fiat-Shamir heuristic; we refer to proofs of this
form as Schnorr proofs. Concretely, $\algo{Prv}$ selects $\alpha, \beta,
\delta, \gamma_0, \ldots, \gamma_{n-1}, \gamma \in \Zp$ and computes the following:
\begin{align*}
  a &= u^\alpha w^\beta \ ; \ b = w^\delta q_n^\gamma
  \prod_{i=0}^{n-1} q_i^{\gamma_i} \ ; \ \algo{chal} = \hash_p(\inst, a, b)
  \\
  \eta &= \alpha + \algo{chal} \varepsilon \ ; \ \nu = \beta + \algo{chal} \iota \ ; \ \xi = \delta +
  \algo{chal} \lambda \ ; \\
  \pi_i &= \gamma_i + \algo{chal} m_i, \ \forall i \in [n] \ ; \ \pi = \gamma +
  \algo{chal} s.
\end{align*}
It then outputs $\Pi = (a, b, \eta, \nu, \xi, \pi_0, \ldots, \pi_{n-1}, \pi)$.
A verifier $\algo{V}$ first checks that $\pairing(u, \tilde{t}) \pairing(x, \ggen) = 1$,
computes $\algo{chal} = \hash_p(\inst, a, b)$, and accepts if and only if the following holds:
\begin{align*}
  a x^\algo{chal} &= u^\eta w^\nu \ ; \ b \hgen^{-\algo{chal}} = w^\xi
  q_n^\pi \prod_{i=0}^{n-1} q_i^{\pi_i}.
\end{align*}

We now show how two provers $\algo{Prv}_0$ and $\algo{Prv}_1$ jointly produce $\Pi$,
where $\algo{Prv}_0$ is a lightweight device that knows $m_{n-1}$ and $\algo{Prv}_1$ is
computationally more powerful and knows $(\varepsilon, s, m_0, \ldots, m_{n-2})$,
following the approach introduced in~\cite{USENIX:HesSinSor23}. Rather than proving
Equation~\eqref{eq:BBSrel}, $\algo{Prv}_0$ and $\algo{Prv}_1$ jointly prove the relation as follows:
\begin{align*}
  \com &= q_{n-1}^{m_{n-1}} q_n^\rho \ \wedge \ \pairing(u, \tilde{t}) \pairing(x,
  \ggen) = 1 \ \wedge \\
  x &= u^\varepsilon w^\iota \ \wedge \ (\com \hgen)^{-1} = w^\lambda q_n^{s'}
  \prod_{i=0}^{n-2} q_i^{m_i}.
\end{align*}
Here, $\inst = (\pk_{\BBS}, u, w, x, \com)$ and $\witness = (\varepsilon, s', \iota, \lambda,
\vec{m}, \rho)$.

\begin{protocolbox}{Joint proof of knowledge of a BBS+ signature}{bbs-joint-proof}

  \begin{description}

    \item[$\Pi \leftarrow \algo{Prv}_0(\gamma_{n-1}, m_{n-1}, \rho)
      \leftrightarrow \algo{Prv}_1(\varepsilon, s, m_0, \ldots, m_{n-2})$ :] The two
      provers jointly compute $\Pi$ as follows.

      \begin{enumerate}

        \item $\algo{Prv}_0$ computes $\com = q_{n-1}^{m_{n-1}} q_n^\rho$ and
          produces a proof $\Pi_0$ that $\inst_0 = (q_{n-1}, q_n, \com)$ and
          $\witness_0 = (m_{n-1}, \rho)$ satisfy $\com = q_{n-1}^{m_{n-1}}
          q_n^\rho$: it selects random $(\gamma_{n-1}, \upsilon)$, computes
          $c = q_{n-1}^{\gamma_{n-1}} q_n^\upsilon$, $\algo{chal}_0 =
          \hash_p(\inst_0, c)$, $\pi_{n-1} = \gamma_{n-1} + \algo{chal}_0
          m_{n-1}$, and $\zeta = \upsilon + \algo{chal}_0 \rho$, and sends
          $\algo{Prv}_1$ the tuple $(\com, \rho, \Pi_0)$ with $\Pi_0 = (c,
          \pi_{n-1}, \zeta)$.

        \item $\algo{Prv}_1$ computes a proof $\Pi_1$ that, for $\inst_1 =
          (\pk_{\BBS}, u, w, x, \com)$ and $\witness_1 = (\varepsilon, s', \iota, \lambda,
          \vec{m}[{:}n-1])$ with $s' = s - \rho$, the following holds:
          \begin{equation}
            \label{eq:reljointproof}
            x = u^\varepsilon w^\iota \ \wedge \ (\com \hgen)^{-1} = w^\lambda
            \prod_{i=0}^{n-2} q_i^{m_i} q_n^{s'}.
          \end{equation}
          To do so, it selects random $(\alpha, \beta)$ and $(\delta,
          \gamma_0, \ldots, \gamma_{n-2}, \gamma)$, computes $a = u^\alpha
          w^\beta$ and $b = w^\delta q_n^\gamma \prod_{i=0}^{n-2}
          q_i^{\gamma_i}$, and derives the following:
          \begin{align*}
            \algo{chal}_1 &= \hash_p(\inst_1, a, b) \ ; \ \eta = \alpha +
            \algo{chal}_1 \varepsilon \ ; \ \nu = \delta + \algo{chal}_1 \iota \ ; \ \xi =
            \delta + \algo{chal}_1 \lambda \ ; \\
            \pi_i &= \gamma_i + \algo{chal}_1 m_i, \ \forall i \in [n-1] \ ; \
            \pi = \gamma + \algo{chal}_1 s'.
          \end{align*}

        \item $\algo{Prv}_1$ outputs $\Pi = (\com, a, b, c, \eta, \nu, \xi,
          \pi_0, \ldots, \pi_{n-1}, \pi, \zeta)$.

      \end{enumerate}

  \end{description}

\end{protocolbox}

\paragraph{Mapping to our solution.}\label{sec:pok-sol} In our protocol, the
secure element $\algo{se}$ and mobile device $\algo{dev}$ jointly produce $\sigma
\leftarrow \NIZKprove(\crs, \inst, \witness)$, where
$\inst = (\gen_0, \gen_1, \pk_{\BBS}, u, w, x, c)$ with $\pk_{\BBS} =
(\ggen, \tilde{t}, q_0, q_1, q_2, \hgen)$, $\witness = (\varepsilon, \iota,
\lambda, s, \algo{uid}, \sk, r)$, and $\rel(\inst, \witness) = 1$ is equivalent to the following:
\begin{align*}
  &\pairing(u, \tilde{t}) \pairing(x, \ggen) = 1 \ \wedge \ x = u^\varepsilon w^\iota \\
  &\wedge \ \hgen^{-1} = w^\lambda q_0^{\algo{uid}} q_1^{\sk} q_2^s \ \wedge \ c =
  \gen_0^{\algo{uid}} \gen_1^r.
\end{align*}
That is, this is an instantiation of Equation~\eqref{eq:BBSrel} with $\vec{m} =
(\algo{uid}, \sk)$ and $c = \gen_0^{\algo{uid}} \gen_1^r$. In this instantiation, $\algo{se}$
executes $\algo{Prv}_0$, outputting $\com = q_1^{\sk} q_2^\rho$, $\rho$, and
proof $\Pi_0$ computed as above except that the challenge $\algo{chal}_0$ also
hashes the message to be signed. $\algo{dev}$ computes a proof $\Pi_1 \leftarrow
\NIZKprove(\crs, \inst_1, \witness_1)$, for $\inst_1
= (\gen_0, \gen_1, \pk_{\BBS}, u, w, x, \com, c)$, $\witness_1 = (\varepsilon, \iota, \lambda,
\algo{uid}, s', r)$, and relation as follows:
\begin{equation*}
  x = u^\varepsilon w^\iota \ \wedge \ (\com \hgen)^{-1} = w^\lambda q_0^{\algo{uid}}
  q_2^{s'} \ \wedge \ c = \gen_0^{\algo{uid}} \gen_1^r.
\end{equation*}
This is computed via sigma protocols, following the same approach as $\Pi_1$ above.

\subsection{Pointcheval-Sanders signatures}\label{sec:ps-sign}

The Pointcheval-Sanders signature
scheme, like BBS+, signs vectors of
messages rather than single messages, using the following
algorithms~\cite{RSA:PoiSan16,RSA:PoiSan18}.

\begin{primitivebox}{Pointcheval-Sanders signatures}{ps}

  \begin{description}

    \item[$(\sk_{\algo{PS}}, \pk_{\algo{PS}}) \leftarrow \algo{PS}\algo{.Gen}(1^\kappa, n)$ :] On
      input security parameter $\kappa$ and integer $n$, randomly select
      $(\chi, \lambda_0, \ldots, \lambda_{n-1}, \iota) \in \Zp^{n+2}$, compute $\tilde{x} =
      \ggen^\chi$, $\tilde{y}_i = \ggen^{\lambda_i}$ for all $i \in [n]$, and
      $\tilde{z} = \ggen^\iota$, and output $\sk_{\algo{PS}} = (\chi, \vec{\lambda}, \iota)$
      and $\pk_{\algo{PS}} = (\tilde{x}, \vec{\tilde{y}}, \tilde{z})$, where
      $\vec{\lambda} = (\lambda_0, \ldots, \lambda_{n-1})$ and $\vec{\tilde{y}} = (\tilde{y}_0,
      \ldots, \tilde{y}_{n-1})$.

    \item[$\psi \leftarrow \algo{PS}\algo{.Sign}(\sk_{\algo{PS}}, \vec{m})$ :] On input
      secret key $\sk_{\algo{PS}}$ and a vector of $n$ messages $\vec{m}$,
      randomly select $k \in \group{1}$ and $\omega \in \Zp$, compute $l = k^{\chi
      + \vec{\lambda} \cdot \vec{m} + \iota\omega}$, and output $\psi = (\omega, k, l)$.

    \item[$\{0,1\} \leftarrow \algo{PS}\algo{.Verify}(\pk_{\algo{PS}}, \vec{m}, \psi)$ :] On
      input public key $\pk_{\algo{PS}}$, vector of messages $\vec{m}$, and
      signature $\psi$, return the result of the check $\pairing(k, \tilde{x}
      \vec{\tilde{y}}^{\vec{m}} \tilde{z}^\omega) = \pairing(l, \ggen)$.

  \end{description}

\end{primitivebox}

The resulting signatures are existentially unforgeable under the
$q$-modified strong Diffie-Hellman ($q$-MSDH) assumption~\cite{RSA:PoiSan18}.

\begin{assumptionbox}{$q$-Modified Strong Diffie-Hellman ($q$-MSDH)}{qmsdh}

  The $q$-MSDH assumption holds if, given $(\gen, \gen^{\chi}, \ggen,
  \ggen^{\chi}, \ggen^{\omega_1}, \ldots, \ggen^{\omega_q}, \gen^{1/(\chi+\omega_1)},
  \ldots, \gen^{1/(\chi+\omega_q)})$ for uniformly random $\chi, \omega_1, \ldots, \omega_q
  \sample \Zp$, no PPT adversary can output a tuple $(\omega, \gen^\nu,
  \gen^{\nu/(\chi+\omega)})$ for $\nu \neq 0$ and $\omega \notin \{\omega_1, \ldots, \omega_q\}$ with
  non-negligible probability.

\end{assumptionbox}

Pointcheval-Sanders signatures can also be extended to support
blind signing: a signer can sign a committed vector without
learning any information about it, cf.~\cite{NDSS:SABMD19}.

\paragraph{Blind signing over committed inputs.}\label{sec:blindsig} Let
$\vec{\gen} = (\gen_0, \ldots, \gen_n) \in \group{1}^{n+1}$ be a vector of
$n+1$ random generators and $\vec{m} = (m_0, \ldots, m_{n-1}) \in \Zp^n$ a
vector of $n$ messages. Signatures over committed
inputs are produced via an interactive
protocol between a signer and a recipient, in which the signer signs a
vector of committed messages such that only the recipient learns
their content~\cite{SCN:CamLys02,ASIACCS:LMPY16}.

The recipient selects an ElGamal secret key $\mu \in \Zp$ with public key
$(\gen, b = \gen^\mu)$, computes $\com = \vec{\gen}^{\vec{m}\|r}$, $k =
\hash_{\group{1}}(\com)$, and $n$ ciphertexts $\ciph_i = (\gen^{\rho_i},
k^{m_i} b^{\rho_i})$ for $i \in [n]$, together with a proof $\Pi \leftarrow
\NIZKprove(\crs, \inst, \witness)$ for $\inst =
(\vec{\gen}, \gen, b, \com, k, (\ciph_i)_{i \in [n]})$, $\witness = ((m_i,
\rho_i)_{i \in [n]}, r)$, and relation as follows:
\begin{align*}
  1 = \vec{\gen}^{\vec{m}\|r} \ \wedge \ \ciph_i = (\gen^{\rho_i},
  k^{m_i}b^{\rho_i}), \ \forall i \in [n].
\end{align*}
It then sends the signer $(\com, \ciph_0, \ldots, \ciph_{n-1}, \gen, b,
\Pi)$. The signer recomputes $k = \hash_{\group{1}}(\com)$, verifies $\Pi$
against $\com$, $k$, and $(\gen, b)$, and rejects if it does not hold.
Otherwise, it parses each $\ciph_i$ as $(e_{i,1}, e_{i,2})$, randomly
selects $\omega \in \Zp$, and computes the following:
\begin{align*}
  \ciph &= (e_1, e_2) = \Bigl(\prod_{i=0}^{n-1} e_{i,1}^{\lambda_i}, \ k^\chi
  \prod_{i=1}^n e_{i,2}^{\lambda_i} k^{\omega\iota}\Bigr).
\end{align*}
It then sends $(\omega, \ciph)$ to the recipient, who decrypts $\ciph = (e_1,
e_2)$ by computing $l = e_2 / e_1^\mu$ and outputs $\psi = (\omega, k, l)$ if it is
a valid Pointcheval-Sanders signature on $\vec{m}$; if the signer followed
the protocol honestly, $l = k^\chi \bigl(\prod_{i=0}^{n-1} k^{\lambda_i m_i}\bigr)
k^{\omega\iota}$.

This protocol adapts readily to partially opened vectors: for $\algo{Disc}
\subseteq [n]$ with $(m_i)_{i \in \algo{Disc}}$ revealed to the signer, the
recipient computes $\com = \prod_{i \in [n]\setminus\algo{Disc}} \gen_i^{m_i}$
and a proof for $\inst = (\algo{Disc}, \vec{\gen}, \gen, b, \com, k,
(\ciph_i)_{i \in [n]\setminus\algo{Disc}})$, $\witness = ((m_i, \rho_i)_{i \in
[n]\setminus\algo{Disc}}, r)$, and relation as follows:
\begin{align*}
  \com &= \gen_n^r \prod_{i \in [n]\setminus\algo{Disc}} \gen_i^{m_i} \ \wedge \
  \ciph_i = (\gen^{\rho_i}, k^{m_i} b^{\rho_i}), \ \forall i \in
  [n]\setminus\algo{Disc}.
\end{align*}
It sends $m_i$ for $i \in \algo{Disc}$ in the clear, and the signer computes the following:
\begin{align*}
  \ciph &= (e_1, e_2) = \Bigl(\prod_{i \in [n]\setminus\algo{Disc}}
  e_{i,1}^{\lambda_i}, \ k^\chi \prod_{i \in \algo{Disc}} k^{m_i \lambda_i} \prod_{i \in
  [n]\setminus\algo{Disc}} e_{i,2}^{\lambda_i} k^{\omega\iota}\Bigr).
\end{align*}
This protocol adapts directly to the withdrawal protocol of
Section~\ref{sec:withdrawal}: we execute it on the vector of messages
$(\algo{uid}_0, v, s_0)$ with randomness $r = 0$.

\paragraph{Proof of withdrawal $\Sigma$.}\label{sec:ps-proof} Let $\tau =
(v, c = \gen_0^{\algo{uid}} \gen_1^s)$ be a token, and let $\pk_{\algo{PS}} = (\tilde{x},
\tilde{y}_0, \tilde{y}_1, \tilde{y}_2, \tilde{z})$ be the public key of the central
bank $\algo{cb}$. Let $\psi = (\omega, k, l)$ be a Pointcheval-Sanders signature that
verifies $1 \gets \algo{PS}\algo{.Verify}(\pk_{\algo{PS}}, (\algo{uid}, v, s), \psi)$, so $\pairing(k,
\tilde{x} \tilde{y}_0^{\algo{uid}} \tilde{y}_1^v \tilde{y}_2^s \tilde{z}^\omega) = \pairing(l,
\ggen)$. To prove that $\tau$ was issued by $\algo{cb}$, the token's owner uses
her mobile device $\algo{dev}$ to compute $\Sigma$ as follows.

\begin{protocolbox}{Proof of withdrawal}{ps-proof}

  \begin{description}

    \item[$\Sigma \leftarrow \algo{ProveWithdrawal}(\pk_{\algo{PS}}, \tau, \omega, \algo{uid}, s)$ :]
      On input the central bank's public key $\pk_{\algo{PS}}$, token $\tau = (v,
      c)$, and witness $(\omega, \algo{uid}, s)$ for a Pointcheval-Sanders signature
      $\psi = (\omega, k, l)$ on $(\algo{uid}, v, s)$, do:

      \begin{enumerate}

        \item Select random $\iota \in \Zp$ and compute $(k^*, l^*) = (k^\iota,
          l^\iota)$.

        \item Generate $\Delta \leftarrow \NIZKprove(\crs, \inst,
          \witness)$, with $\inst = (\gen_0, \gen_1, \tilde{x}, \tilde{y}_0,
          \tilde{y}_1, \tilde{y}_2, \tilde{z}, v, c, k^*, l^*)$, $\witness = (\omega, \algo{uid}, s)$,
          and relation as follows:
          \begin{align*}
            \pairing(l^*, \ggen) &= \pairing(k^*, \tilde{x} \tilde{y}_0^{\algo{uid}} \tilde{y}_1^v
            \tilde{y}_2^s \tilde{z}^\omega) \ \wedge \ c = \gen_0^{\algo{uid}} \gen_1^s.
          \end{align*}

        \item Return $\Sigma = (k^*, l^*, \Delta)$.

      \end{enumerate}

  \end{description}

\end{protocolbox}

Since $s_0$ is random and not known, $(k^*, l^*)$ leaks no information about
$\algo{uid}_0$. To compute $\Delta$, $\algo{dev}$ first computes $\tilde{c} =
\tilde{y}_0^{\algo{uid}} \tilde{y}_2^s \tilde{z}^\omega$, randomly selects $\alpha, \beta,
\delta \in \Zp$, computes $a = \gen_0^\alpha \gen_1^\beta$, $\tilde{a} =
\tilde{y}_0^\alpha \tilde{y}_2^\beta \tilde{z}^\delta$, and challenge $\gamma =
\hash_p(\inst, \tilde{c}, a, \tilde{a})$, then computes $\pi_{\algo{uid}} =
\alpha + \gamma \algo{uid}$, $\pi_s = \beta + \gamma s$, and $\pi_\omega = \delta +
\gamma \omega$, and sets $\Delta = (\tilde{c}, a, \tilde{a}, \pi_{\algo{uid}}, \pi_s,
\pi_\omega)$. These are standard Schnorr proofs for $\inst' = (\gen_0,
\gen_1, \tilde{y}_0, \tilde{y}_2, \tilde{z}, c, \tilde{c})$, $\witness' = (\omega, \algo{uid},
s)$, and $\rel(\inst', \witness') = 1$ equivalent to $c = \gen_0^{\algo{uid}}
\gen_1^s \wedge \tilde{c} = \tilde{y}_0^{\algo{uid}} \tilde{y}_2^s \tilde{z}^\omega$.

If $(\omega, k^*, l^*)$ is a valid Pointcheval-Sanders signature on $(\algo{uid}, v,
s)$ under $\pk_{\algo{PS}}$, then $\pairing(l^*, \ggen) = \pairing(k^*, \tilde{x} \tilde{c}
\tilde{y}_1^v)$, so a verifier, given $\Sigma = (k^*, l^*, \Delta)$ and token
$\tau = (v, c)$, checks that $\tau$ was issued by $\algo{cb}$ by parsing $\Delta =
(\tilde{c}, a, \tilde{a}, \pi_{\algo{uid}}, \pi_s, \pi_\omega)$, computing $\gamma =
\hash_p(\inst', a, \tilde{a})$, and verifying the following:
\begin{align*}
  \pairing(l^*, \ggen) &= \pairing(k^*, \tilde{x} \tilde{c} \tilde{y}_1^v) \ \wedge \
  c^\gamma a = \gen_0^{\pi_{\algo{uid}}} \gen_1^{\pi_s} \ \wedge \
  \tilde{c}^\gamma \tilde{a} = \tilde{y}_0^{\pi_{\algo{uid}}} \tilde{y}_2^{\pi_s}
  \tilde{z}^{\pi_\omega}.
\end{align*}
